\subsection{Path Discovery con RIPE Atlas}\label{sec:ripe}


\subsubsection{Valparaíso - IP y ASN por salto}


\begingroup
\fontsize{9}{11}\selectfont
\begin{longtable}{p{0.95cm}p{4.716666666666667cm}p{4.716666666666667cm}p{4.716666666666667cm}}
\caption{RIPE de ida: IP y ASN por TTL, Valparaíso, ronda 2.}\label{ripe:138.204.230.242} \\
\toprule
\textbf{HOP} & \textbf{UDP} & \textbf{ICMP} & \textbf{TCP} \\ \midrule
\endfirsthead
\toprule
\textbf{HOP} & \textbf{UDP} & \textbf{ICMP} & \textbf{TCP} \\ \midrule
\endhead
1 & 192.168.100.1 [priv./esp.] & 192.168.100.1 [priv./esp.] & 192.168.100.1 [priv./esp.] \\
2 & 100.88.0.1 [priv./esp.] & 100.88.0.1 [priv./esp.] & 100.88.0.1 [priv./esp.] \\
3 & 10.87.0.3 [priv./esp.] & 10.87.0.2 [priv./esp.] & 10.87.0.3 [priv./esp.] \\
4 & 10.87.0.46 [priv./esp.] & 10.87.0.46 [priv./esp.] & 10.87.0.46 [priv./esp.] \\
5 & 45.68.16.129 [AS n/d] & 45.68.16.129 [AS n/d] & 45.68.16.129 [AS n/d] \\
6 & 10.100.154.1 [priv./esp.] & 10.100.154.1 [priv./esp.] & 10.100.154.1 [priv./esp.] \\
7 & 138.204.228.164 [AS61506] & 138.204.228.164 [AS61506] & 138.204.228.164 [AS61506] \\
8 & 138.204.228.162 [AS61506] & 138.204.228.162 [AS61506] & 138.204.228.162 [AS61506] \\
9 & 138.204.230.242 [AS61506] & 138.204.230.242 [AS61506] & 138.204.230.242 [AS61506] \\
\bottomrule
\end{longtable}
\endgroup
\begingroup
\fontsize{9}{11}\selectfont
\addtocounter{table}{-1}
\begin{longtable}{p{3.1cm}p{1.9cm}p{2.1cm}p{2.1cm}p{1.8cm}p{2.1cm}}
\toprule
\textbf{Método} & \textbf{IP únicas} & \textbf{Cobertura} & \textbf{Estabilidad} & \textbf{Sondas} & \textbf{Llegada / 3} \\ \midrule
\endfirsthead
\toprule
\textbf{Método} & \textbf{IP únicas} & \textbf{Cobertura} & \textbf{Estabilidad} & \textbf{Sondas} & \textbf{Llegada / 3} \\ \midrule
\endhead
ICMP & 9 & 100.0\% & 1.000 & 27 & 3 \\
TCP & 9 & 100.0\% & 1.000 & 27 & 3 \\
UDP & 9 & 100.0\% & 1.000 & 27 & 3 \\
\bottomrule
\end{longtable}
\endgroup

En Valparaíso, RIPE y Scamper cuentan la misma historia, nueve IP únicas, cobertura de 100\% estabilidad de 1.000 y llegada al destino en las tres rondas con los tres protocolos (ICMP, TCP y UDP) quedan empatados como los más informativos). La diferencia es el costo de 27 sondas contra 9,porque RIPE envía por defecto varios paquetes por salto. También se repite el cambio en el TTL 3, con 10.87.0.2 para ICMP y 10.87.0.3 para UDP y TCP, igual que en Scamper. Dos herramientas distintas coinciden entonces en que, en ese punto de la red del operador, ICMP y UDP/TCP no siguen la misma interfaz.


\subsubsection{La Habana - IP y ASN por salto}
\begingroup
\fontsize{9}{11}\selectfont
\begin{longtable}{p{0.95cm}p{4.716666666666667cm}p{4.716666666666667cm}p{4.716666666666667cm}}
\caption{RIPE de ida: IP y ASN por TTL, La Habana, ronda 2.}\label{ripe:200.14.48.134} \\
\toprule
\textbf{HOP} & \textbf{UDP} & \textbf{ICMP} & \textbf{TCP} \\ \midrule
\endfirsthead
\toprule
\textbf{HOP} & \textbf{UDP} & \textbf{ICMP} & \textbf{TCP} \\ \midrule
\endhead
1 & 192.168.100.1 [priv./esp.] & 192.168.100.1 [priv./esp.] & 192.168.100.1 [priv./esp.] \\
2 & 100.88.0.1 [priv./esp.] & 100.88.0.1 [priv./esp.] & 100.88.0.1 [priv./esp.] \\
3 & 10.87.0.3 [priv./esp.] & 10.87.0.2 [priv./esp.] & 10.87.0.3 [priv./esp.] \\
4 & 10.87.0.46 [priv./esp.] & 10.87.0.46 [priv./esp.] & 10.87.0.46 [priv./esp.] \\
5 & * & * & * \\
6 & 154.54.92.41 [AS174] & 154.54.92.41 [AS174] & 154.54.92.41 [AS174] \\
7 & *; 154.54.29.157 [AS174] & 154.54.29.157 [AS174] & *; 154.54.29.157 [AS174]; 154.54.3.29 [AS174] \\
8 & 154.54.167.185 [AS174] & 154.54.167.185 [AS174] & 154.54.167.185 [AS174] \\
9 & 38.88.57.26 [AS174] & 38.104.94.178 [AS174] & 38.88.57.26 [AS174] \\
10 & * & * & * \\
11 & 69.79.104.1 [AS23520] & 190.242.168.37 [AS23520] & *; 190.242.168.39 [AS23520] \\
12 & * & 200.0.16.21 [AS11960] & 200.0.16.16 [AS11960] \\
13 & 190.242.168.39 [AS23520] & 200.0.16.29 [AS11960] & 200.0.16.21 [AS11960] \\
14 & 200.0.16.16 [AS11960] & * & 200.0.16.25 [AS11960] \\
15 & 200.0.16.21 [AS11960] & * & *; 200.0.16.25 [AS11960] \\
16 & 200.0.16.29 [AS11960] & 200.14.48.178 [AS27725] & *; 192.168.100.1 [priv./esp.] \\
17 & * & 200.14.48.134 [AS27725] & *; 100.88.0.1 [priv./esp.]; 200.0.16.29 [AS11960] \\
18 & * & -- & * \\
19 & 200.14.48.178 [AS27725] & -- & 200.14.48.178 [AS27725] \\
20 & 200.14.48.134 [AS27725] & -- & 200.14.48.134 [AS27725] \\
\bottomrule
\end{longtable}
\endgroup
\begingroup
\fontsize{9}{11}\selectfont
\addtocounter{table}{-1}
\begin{longtable}{p{3.1cm}p{1.9cm}p{2.1cm}p{2.1cm}p{1.8cm}p{2.1cm}}
\toprule
\textbf{Método} & \textbf{IP únicas} & \textbf{Cobertura} & \textbf{Estabilidad} & \textbf{Sondas} & \textbf{Llegada / 3} \\ \midrule
\endfirsthead
\toprule
\textbf{Método} & \textbf{IP únicas} & \textbf{Cobertura} & \textbf{Estabilidad} & \textbf{Sondas} & \textbf{Llegada / 3} \\ \midrule
\endhead
TCP & 14 & 82.4\% & 0.730 & 57 & 3 \\
UDP & 14 & 77.8\% & 0.753 & 54 & 3 \\
ICMP & 13 & 76.5\% & 1.000 & 51 & 3 \\
\bottomrule
\end{longtable}
\endgroup

%En La Halana, RIPE llega al destino en las tres rondas con los tres protocolos, pero con trayectos de largo distinto, con ICMP responde en el TTL 17, igual que Scamper con icmp, y con UDP y TCP en el TTL 20. En la misma ronda, las sondas UDP y TCP necesitan tres saltos más que las ICMP y atraviesan más interfaces de A11969, lo que coincide con lo visto en Scamper. El protocolo más informativo es TCP, con 14 IP únicas y cobertura de 82.4\%, seguido de UDP (14 IP, 77,8\%) e ICMP (13 IP, 76,5\%).

Para Habana, las pruebas muestran que ICMP alcanza 17 saltos, mientras que UDP y TCP llegan a 20, atravesando más interfaces de AS11960. En este caso, TCP ofrece la mayor cobertura, con 14 UP únicas y un 82,4\% seguido de UDP (14 IP y 77,8\%) e ICMP (13 IP y 76,5\%). Sin embargo, ICMP es más estable entre rondas 1.000 y requiere menos sondas (51) frente a 57 de TCP, 73,0\% cobertura  y UDP con 54 sondas (75,3\%). Las respuestas múltiples por TTL sugieren balanceo de carga, mientras que las direcciones privadas en los TTL altos se conservan como registros de RIPE, pero se excluyen de la interpretación de la ruta.

%Protocolo elegido: TCP.

\subsubsection{Rusia - IP y ASN por salto}
\begingroup
\fontsize{9}{11}\selectfont
\begin{longtable}{p{0.95cm}p{4.716666666666667cm}p{4.716666666666667cm}p{4.716666666666667cm}}
\caption{RIPE de ida: IP y ASN por TTL, Rusia, ronda 2.}\label{ripe:37.79.226.46} \\
\toprule
\textbf{HOP} & \textbf{UDP} & \textbf{ICMP} & \textbf{TCP} \\ \midrule
\endfirsthead
\toprule
\textbf{HOP} & \textbf{UDP} & \textbf{ICMP} & \textbf{TCP} \\ \midrule
\endhead
1 & 192.168.100.1 [priv./esp.] & 192.168.100.1 [priv./esp.] & 192.168.100.1 [priv./esp.] \\
2 & 100.88.0.1 [priv./esp.] & 100.88.0.1 [priv./esp.] & 100.88.0.1 [priv./esp.] \\
3 & 10.87.0.3 [priv./esp.] & 10.87.0.2 [priv./esp.] & 10.87.0.3 [priv./esp.] \\
4 & 10.87.0.46 [priv./esp.] & 10.87.0.46 [priv./esp.] & 10.87.0.46 [priv./esp.] \\
5 & * & * & * \\
6 & 154.54.92.41 [AS174] & 154.54.92.41 [AS174] & *; 154.54.92.41 [AS174] \\
7 & *; 154.54.29.157 [AS174] & *; 154.54.3.29 [AS174] & *; 154.54.3.29 [AS174] \\
8 & 154.54.167.185 [AS174] & 154.54.168.85 [AS174] & 154.54.168.85 [AS174] \\
9 & 154.54.10.210 [AS174] & 154.54.10.210 [AS174] & 154.54.10.210 [AS174] \\
10 & 195.22.211.61 [AS6762] & 195.22.211.61 [AS6762] & 195.22.211.209 [AS6762] \\
11 & *; 100.88.0.1 [priv./esp.]; 192.168.100.1 [priv./esp.]; 195.22.214.161 [AS6762] & * & * \\
12 & * & * & * \\
13 & * & * & * \\
14 & 2.63.212.246 [AS12389] & 2.63.212.246 [AS12389] & *; 2.63.212.246 [AS12389] \\
15 & * & 37.79.226.46 [AS12389] & 37.79.226.46 [AS12389] \\
16 & * & -- & 37.79.226.46 [AS12389] \\
17 & * & -- & *; 37.79.226.46 [AS12389] \\
18 & * & -- & -- \\
19 & * & -- & -- \\
255 & * & -- & -- \\
\bottomrule
\end{longtable}
\endgroup
\begingroup
\fontsize{9}{11}\selectfont
\addtocounter{table}{-1}
\begin{longtable}{p{3.1cm}p{1.9cm}p{2.1cm}p{2.1cm}p{1.8cm}p{2.1cm}}
\toprule
\textbf{Método} & \textbf{IP únicas} & \textbf{Cobertura} & \textbf{Estabilidad} & \textbf{Sondas} & \textbf{Llegada / 3} \\ \midrule
\endfirsthead
\toprule
\textbf{Método} & \textbf{IP únicas} & \textbf{Cobertura} & \textbf{Estabilidad} & \textbf{Sondas} & \textbf{Llegada / 3} \\ \midrule
\endhead
TCP & 11 & 82.4\% & 0.840 & 51 & 3 \\
ICMP & 11 & 73.3\% & 0.889 & 45 & 3 \\
UDP & 11 & 55.0\% & 0.714 & 60 & 0 \\
\bottomrule
\end{longtable}
\endgroup

%Protocolo(s) elegido(s): TCP. 

En Rusia, el resultado de RIPE depende mucho del protocolo. Con TCP e ICMP llega al destino en las tres rondas, TCP obtiene 11 IP únicas, cobertura de 82.4\%, estabilidad de 0.840 y 51 sondas, e ICMP también 11 IP, con 73.3\%, 0.889 y 45 sondas. Con UDP no llega nunca (0 de 3): ve las mismas 11 IP pero con cobertura de 55.0\%, estabilidad de 0.714 y más sondas (60). Es el mismo comportamiento que mostraron udp y udp-paris en Scamper, de modo que dos herramientas con implementaciones distintas coinciden en que el destino ruso no responde a las sondas UDP. Comparado con Scamper, RIPE no mejora lo descubierto (11 IP contra 12 de icmp) y solo mejora levemente la cobertura con TCP (82.4 \% contra 80.0 \%), a tres veces el costo. En el TTL 11 de UDP vuelven a aparecer 192.168.100.1 y 100.88.0.1, que trato igual que arriba.

\Needspace{9\baselineskip}
\subsubsection{ripe / traceroute / UDP}
Valparaíso como ejemplo. Entrenamiento: ripe\_\allowbreak{}20261026T205835\_\allowbreak{}d73ba97e; contraste: rondas 2 y 3. Las reglas completas de los tres destinos están en el archivo de predicciones enlazado en el anexo.

\begingroup
\fontsize{9}{11}\selectfont
\addtocounter{table}{-1}
\begin{longtable}{p{3.1cm}p{5.8cm}p{5.1cm}}
\toprule
\textbf{Campo} & \textbf{Regla en ronda 1} & \textbf{Contraste posterior} \\ \midrule
\endfirsthead
\toprule
\textbf{Campo} & \textbf{Regla en ronda 1} & \textbf{Contraste posterior} \\ \midrule
\endhead
ip.ttl & incremento 1 cada 3 paquetes & r2: cumple; r3: cumple \\
ip.id & variable en la captura & r2: no exacto; r3: no exacto \\
ip.len & constante 76 & r2: cumple; r3: cumple \\
ip.proto & constante 17 & r2: cumple; r3: cumple \\
udp.srcport & constante 20480 & r2: cumple; r3: cumple \\
udp.dstport & constante 33435 & r2: cumple; r3: cumple \\
udp.length & constante 56 & r2: cumple; r3: cumple \\
udp.checksum & incremento 1 por paquete & r2: cumple; r3: cumple \\
\bottomrule
\end{longtable}
\endgroup
Payload: [48] bytes. Hexadecimal inicial (offset 0):

\begin{quote}\ttfamily\footnotesize\raggedright 0352 6874 7470 3a2f 2f61 746c 6173 2e72 6970 652e 6e65 7420 4174 6c61 7320 7361 7973 2048 6921 0000 0000 0000 0000 0000\end{quote}
Texto ASCII desde offset 2: http://atlas.ripe.net Atlas says Hi!

Al sumar payload\_\allowbreak{}word0\_\allowbreak{}BE, udp.checksum, udp.srcport, udp.dstport módulo 65535, el resultado se mantiene constante durante la ejecución. Ese valor base se aprende del primer paquete de cada ejecución y no lo doy por universal. Es lo que se esperaría si el emisor compensara el checksum con el payload.

Predicción concreta en ronda 2, campo ip.ttl: la primera solicitud tiene valor 1. Aplicando la regla aprendida (y esa base, cuando corresponde) predije los valores de las solicitudes 2 a 4 antes de compararlos con lo que realmente se observó.

\begingroup
\fontsize{9}{11}\selectfont
\addtocounter{table}{-1}
\begin{longtable}{p{3.2cm}p{5.1cm}p{5.1cm}}
\toprule
\textbf{Solicitud / trama} & \textbf{Predicción} & \textbf{Observación} \\ \midrule
\endfirsthead
\toprule
\textbf{Solicitud / trama} & \textbf{Predicción} & \textbf{Observación} \\ \midrule
\endhead
2 / 7 & 1 & 1 \\
3 / 10 & 1 & 1 \\
4 / 12 & 2 & 2 \\
\bottomrule
\end{longtable}
\endgroup
Alcance: 27 solicitudes en el ejemplo inicial. Una verificación de campo o byte no es igual a una medición independiente, y a los campos que varían les dejo sus restricciones en lugar de inventarles un valor.

\Needspace{9\baselineskip}
\subsubsection{ripe / traceroute / ICMP}
Ejemplo: Valparaíso. Entrenamiento: ripe\_\allowbreak{}20261026T210015\_\allowbreak{}4d0ba1ee, con contraste de rondas 2 y 3.

\begingroup
\fontsize{9}{11}\selectfont
\addtocounter{table}{-1}
\begin{longtable}{p{3.1cm}p{5.8cm}p{5.1cm}}
\toprule
\textbf{Campo} & \textbf{Regla en ronda 1} & \textbf{Contraste posterior} \\ \midrule
\endfirsthead
\toprule
\textbf{Campo} & \textbf{Regla en ronda 1} & \textbf{Contraste posterior} \\ \midrule
\endhead
ip.ttl & incremento 1 cada 3 paquetes & r2: cumple; r3: cumple \\
ip.id & variable en la captura & r2: no exacto; r3: no exacto \\
ip.len & constante 76 & r2: cumple; r3: cumple \\
ip.proto & constante 1 & r2: cumple; r3: cumple \\
icmp.type & constante 8 & r2: cumple; r3: cumple \\
icmp.code & constante 0 & r2: cumple; r3: cumple \\
icmp.ident & constante 36864 & r2: cumple; r3: cumple \\
icmp.seq & incremento 1 por paquete & r2: cumple; r3: cumple \\
icmp.checksum & constante 0x0001 & r2: cumple; r3: cumple \\
\bottomrule
\end{longtable}
\endgroup
Payload: [48] bytes. Hexadecimal inicial (offset 0):

\begin{quote}\ttfamily\footnotesize\raggedright d4dc 6874 7470 3a2f 2f61 746c 6173 2e72 6970 652e 6e65 7420 4174 6c61 7320 7361 7973 2048 6921 0000 0000 0000 0000 0000\end{quote}
Texto ASCII desde offset 2: http://atlas.ripe.net Atlas says Hi!

Sumando payload\_\allowbreak{}word0\_\allowbreak{}BE, icmp.checksum, icmp.ident, icmp.seq módulo 65535 se mantuvo constante durante la ejecución. Ese valor base se aprende del primer paquete de cada ejecución y no lo doy por universal

Predicción concreta en ronda 2, campo icmp.seq: la primera solicitud tiene valor 1. Con la regla aprendida y esa base cuando corresponde, se predicen las solicitudes 2 a 4 antes de contrastarlas con sus valores observados.

\begingroup
\fontsize{9}{11}\selectfont
\addtocounter{table}{-1}
\begin{longtable}{p{3.2cm}p{5.1cm}p{5.1cm}}
\toprule
\textbf{Solicitud / trama} & \textbf{Predicción} & \textbf{Observación} \\ \midrule
\endfirsthead
\toprule
\textbf{Solicitud / trama} & \textbf{Predicción} & \textbf{Observación} \\ \midrule
\endhead
2 / 3 & 2 & 2 \\
3 / 5 & 3 & 3 \\
4 / 7 & 4 & 4 \\
\bottomrule
\end{longtable}
\endgroup
Alcance: 27 solicitudes en el ejemplo inicial.

\Needspace{9\baselineskip}
\subsubsection{ripe / traceroute / TCP}
Valparaíso como ejemplo. Entrenamiento: ripe\_\allowbreak{}20261026T210108\_\allowbreak{}91a49ffa; contraste: rondas 2 y 3.

\begingroup
\fontsize{9}{11}\selectfont
\addtocounter{table}{-1}
\begin{longtable}{p{3.1cm}p{5.8cm}p{5.1cm}}
\toprule
\textbf{Campo} & \textbf{Regla en ronda 1} & \textbf{Contraste posterior} \\ \midrule
\endfirsthead
\toprule
\textbf{Campo} & \textbf{Regla en ronda 1} & \textbf{Contraste posterior} \\ \midrule
\endhead
ip.ttl & incremento 1 cada 3 paquetes & r2: cumple; r3: cumple \\
ip.id & variable en la captura & r2: no exacto; r3: no exacto \\
ip.len & constante 88 & r2: cumple; r3: cumple \\
ip.proto & constante 6 & r2: cumple; r3: cumple \\
tcp.srcport & constante 20480 & r2: cumple; r3: cumple \\
tcp.dstport & constante 80 & r2: cumple; r3: cumple \\
tcp.flags & constante 0x0002 & r2: cumple; r3: cumple \\
tcp.seq\_\allowbreak{}raw & incremento 1 por paquete & r2: cumple; r3: cumple \\
tcp.ack\_\allowbreak{}raw & constante 0 & r2: cumple; r3: cumple \\
tcp.len & constante 48 & r2: cumple; r3: cumple \\
tcp.checksum & incremento -1 por paquete & r2: cumple; r3: cumple \\
\bottomrule
\end{longtable}
\endgroup
Payload: [48] bytes. Hexadecimal inicial (offset 0):

\begin{quote}\ttfamily\footnotesize\raggedright 6874 7470 3a2f 2f61 746c 6173 2e72 6970 652e 6e65 7420 4174 6c61 7320 7361 7973 2048 6921 0000 0000 0000 0000 0000 0000\end{quote}
Texto ASCII desde offset 0: http://atlas.ripe.net Atlas says Hi!

Predicción concreta en ronda 2, campo tcp.seq\_\allowbreak{}raw: la primera solicitud tiene valor 1.

\begingroup
\fontsize{9}{11}\selectfont
\addtocounter{table}{-1}
\begin{longtable}{p{3.2cm}p{5.1cm}p{5.1cm}}
\toprule
\textbf{Solicitud / trama} & \textbf{Predicción} & \textbf{Observación} \\ \midrule
\endfirsthead
\toprule
\textbf{Solicitud / trama} & \textbf{Predicción} & \textbf{Observación} \\ \midrule
\endhead
2 / 3 & 2 & 2 \\
3 / 5 & 3 & 3 \\
4 / 7 & 4 & 4 \\
\bottomrule
\end{longtable}
\endgroup

Alcance: 27 solicitudes del ejemplo inicial. Verificar un campo o un byte no es lo mismo que una medición independiente, y a los campos que varían les dejo sus restricciones en lugar de inventarles un valor.

\subsubsection{Lectura conjunta de los patrones de RIPE}

RIPE es la herramienta más fácil de reconocer en una captura, porque sus paquetes llevan un texto legible, en los tres protocolos el payload tiene 48 bytes y contiene \verb| http://atlas.ripe.net Atlas says Hi!|, a partir  del offset 2 en UDP e ICMP (los dos primeros bytes varían, 0352 y d4dc en los ejemplos) y desde el offset 0 en TCP. Scamper, en cambio no incluye texto alguno en sus payloads. Las reglas de campos también se cumplen, ip.ttl parte en 1 en UDP, icmp.seq en 1 en ICMP y tcp.seq\_saw en 1 en TCP. En el ejemplo, cada ficha reúne 27 solicitudes, coherente con tres paquetes por salto.


